\section{Reader, Planner, and Repair Interfaces}
\label{app:prompts}

The prompt builders and provider schemas are generated by
\path{src/procurement_graph/reasoning/typed_planning.py}.  Its audited SHA256 is recorded in
\path{paper/tmlr_v3/artifacts/REPRODUCIBILITY_MANIFEST.md}, together with the runner and
configuration hashes.  The historical teacher summary preserves both model names but not a full
command or prompt version.  Table \ref{tab:teacher-call-contract} therefore distinguishes stored
facts from settings reconstructed from the audited runner.  The code is the canonical source for
exact whitespace, dynamic schema context, and enum expansion.

\begin{table}[h]
  \centering
  \caption{Teacher prompt and call contract.  ``Stored'' is present in the harvest summary;
  ``reconstructed'' is implied by the audited runner and model-dependent prompt selection.}
  \label{tab:teacher-call-contract}
  \small
  \begin{tabularx}{\linewidth}{p{0.18\linewidth}p{0.19\linewidth}p{0.22\linewidth}X}
    \toprule
    Call & Model & Output control & Sampling and visibility \\
    \midrule
    Step-1 intent & \texttt{gpt-5.4-nano} (stored) & strict intent-program JSON Schema & temperature 0; question and schema slice; no records or oracle \\
    Step-2 graph & \texttt{grok-4-1-fast-}\newline\texttt{non-reasoning} (stored) & reconstructed \texttt{lean}/\texttt{optional} strict graph schema & temperature 0; at most two structural samples; question, Step-1 output, and schema slice \\
    Repair reading & Step-1 model & fixed eight-section text scaffold & called for semantic repair; hidden-answer fields removed \\
    Repair graph & Step-2 model & strict repair/graph JSON Schema & at most two teacher-harvest repairs; failed graph and structured feedback, but no oracle value \\
    \bottomrule
  \end{tabularx}
\end{table}

The shared API client uses a 60-second timeout and allows three retries after the initial request,
with waits of 2, 4, and 6 seconds.  Rate limits remain retryable; most other 4xx failures terminate
immediately.  The runner does not set an explicit output-token limit.  These call settings are
reconstructed from the current client because they were not copied into the historical summary.

\subsection{Step-1 understander}

The exact system instruction is:

\begin{quote}\small
``You are the typed intent programmer for a UK public-procurement KGQA system. Convert the question
into a small computation program. Do not answer the question. Use only information explicitly
present in the question. The output shape is enforced by a strict JSON schema.''
\end{quote}

The user JSON contains the keys \texttt{question}, \texttt{task}, \texttt{allowed\_slots},
\texttt{operation\_guidance}, \texttt{hard\_rules}, and
\texttt{retrieved\_schema\_context}.  The required output contains:

\begin{verbatim}
answer_signature: operation, value_type, answer_field_text
program: ordered typed steps with literal filters and dependencies
answer_step: id of the final step
unsupported_or_ambiguous: explicit reason or empty
\end{verbatim}

Allowed literal slots are publication year/range, CPV code and label, procurement category limited
to goods/services/works, buyer, supplier, exact quoted title, date, and numeric amount.  Root steps
cannot invent an ``all records'' input.  Later steps may reference only earlier identifiers.
Bridge intents must construct source records, a distinct entity set, a dependent target set, and a
final operation.  Future years remain executable no-result queries rather than predicted
abstentions.  If the question is unsupported or ambiguous, the official program is empty.

\subsection{Step-2 graph planner}

Step 2 is called only when the Step-1 program cannot take the deterministic fast path or is vetoed.
The teacher model name selects the lean prompt rather than the larger capability-card prompt.  Its
exact system instruction is:

\begin{quote}\small
``You are the structured graph planner for a UK public-procurement KGQA system. Turn the question
and the Step-1 briefing into an executable \texttt{graph\_plan} (variables + return). A strict
schema fixes the output shape; fill the fields faithfully and never answer the question.''
\end{quote}

It receives the original question, the Step-1 briefing, schema context, and instruction list.
Its JSON graph contains a flat variable array, relations, a return record, a template label, and an
optional abstention reason.  Record filters contain only literal constraints from the question.
Derived phrases are represented as variables and relations, not copied into filter slots.  A
bridge must emit a typed entity set and bind it into the target record variable.  Money sums must
request the additive-value guard.  Unused schema fields take declared empty values rather than
invented semantics.

The planner never receives records, an execution result, or the oracle answer.  The consistency
checker compares the graph with explicit years, CPV codes, categories, titles, numeric literals,
buyer/supplier roles, and the requested count target.  A structural resample appends only the
diagnostic that compilation or consistency failed and asks for a different, simpler decomposition.

\subsection{Feedback-conditioned repair}

An eligible runtime failure produces a structured record containing its stage, code, path, and
message.  Repair reading is instructed to revise the understanding from the failed graph and
feedback without answering or inferring a reference answer.  Graph repair is instructed to
diagnose the failed graph, choose an allowed action, and emit a graph under the original schema.
The feedback sanitizer removes gold/reference/expected/oracle-answer keys.  Repair SFT conditions
on the original question, Step-1 briefing, failed graph, and this visible feedback.  It targets the
accepted repaired graph.  Oracle mismatch feedback states only that external validation failed and
omits the expected value.  Clean no-result and unsupported cases terminate; they are not repeatedly
replanned in search of a non-empty answer.  The teacher collection budget is two repairs, while the
260-question evaluation budget is one.

\subsection{Training examples and information boundaries}

Direct-plan SFT serialises $(x,u)\mapsto z^+$.  Repair SFT serialises
$(x,u,z^-,f)\mapsto z^+$.  Preference training uses the same question and briefing with accepted
and rejected graph completions.  Neither the oracle answer nor the candidate's returned value is
included in the student input or target.  This boundary is important for interpreting the learned
signal: the model can imitate the selected planning decisions and the response to visible failure
feedback, but it is never directly told which individual decision made the denotation correct.

Offline selection nevertheless uses more than runtime grounding.  A direct answerable target must
pass execution/release checks, match the hidden answer oracle, and have the expected output shape.
A runtime-passing but answer-wrong graph enters the hard-negative path.  Its repair becomes a
repair-SFT or DPO positive only after deterministic re-execution matches the same hidden gates.
Thus grounding defines an executable boundary, whereas the oracle filters denotational correctness;
neither proves that every question constraint is present in the plan.

The teacher writer is append-only and resume is not transactional.  It records the trace before
the downstream SFT/DPO sinks, while resume skips every identifier already in the trace file.  A
failure between those writes can therefore leave an incomplete downstream corpus.  The emitted
metadata string \texttt{acceptance: executor\_verifier} is also stale: the actual control flow has
already applied the oracle/shape branch before direct SFT is written.  Reproduction should use the
runner logic and per-sink counts rather than that label.
